\begin{textsample}{NEG Dim 6 – vem\_ed\_14 – Score 3.00 – vem\_ed\_14/t060.txt}  \label{ex:f6_neg_080}
PCIs EM CONGRESSOS INTERNACIONAIS E PUBLICAÇÕES Os Projetos Colaborativos Internacionais (PCIs / Cesu) têm marcado presença em congressos internacionais e publicações acadêmicas.
Osvaldo Succi Junior, coordenador da equipe dos Projetos Colaborativos Internacionais (PCI / Cesu / CPS), assina dois capítulos do livro The Guide to COIL Virtual Exchange.
O volume foi lançado em 26 de outubro de 2022 na principal conferência mundial sobre o tema: IVEC (International Virtual Exchange Conference), em Valencia (Espanha).
Está disponível para compra em: https://bit.ly/3hc4vZq.
Os organizadores do livro são Jon Rubin e Sarah Guth, referências mundiais em COIL (acrônimo para Collaborative Online International Learning, ou aprendizagem internacional colaborativa online).
Um dos capítulos com a assinatura de Succi Junior aborda projetos de intercâmbio virtual / COIL ao redor do mundo e outro, a escalabilidade desses projetos.
O PCI realizado em 2019 pelos professores Camila Kami (Fatec Bauru), Luiz Roberto Madureira Iório e Liesa Kyer (BridgeValley Community College, EUA) é mencionado no livro.
Além do lançamento do livro, Succi Junior participou de um painel em 27 de outubro, durante a conferência IVEC.
Sua apresentação abordou “Desafios e oportunidades culturais nas Américas: uma perspectiva brasileira”.
Ricardo Nóbrega, coordenador do curso de Comércio Exterior da Fatec Indaiatuba, também teve trabalho aprovado, que foi apresentado por sua colega Mona Pearl (DePaul University, EUA).
Os professores são parceiros desde 2019 em PCIs, relatados no artigo.
Patrícia Sales Patrício, professora da Fatec Ipiranga e responsável pela comunicação dos PCIs / Cesu, fez a apresentação “Em busca do signo da \textbf{relação}: relato de experiência de um intercâmbio virtual Brasil – Portugal”, no 1º Congresso Internacional de Interculturalidade, Lusofonia, Identidade e Comunicação (ILIC\_22).
Patrícia desenvolve, desde novembro de 2021, um PCI com a professora Ana Torres, da Universidade de Aveiro (Portugal).
No primeiro semestre de 2022, a professora portuguesa Paula Pinheiro se uniu ao projeto.
Maira Rezende (Fatec Sorocaba) também apresentou no ILIC\_22 o trabalho “Interculturalidade no Desenvolvimento de Projetos Colaborativos entre Instituições de Ensino” com seu colega de PCI na Inacap (Chile), Patricio Hernán Marabolí Albornoz.
Os autores publicaram sua experiência em capítulo do e-book Ensaios: Educação, Ciência e Tecnologia.

% matched lemmas: relação
\end{textsample}
